\documentclass[tikz, border=2mm]{standalone}
\usepackage{tikz}
\usepackage{amsmath} % For \frac
\usepackage{amssymb} % For symbols if needed
\usetikzlibrary{math} % For math functions like abs()

\begin{document}

\begin{tikzpicture}[scale=1.2] % Adjust scale as needed

% Define k (must be non-zero, using a positive value for visualization)
\def\k{-1}

% Z-plane
\begin{scope}
    \draw[->] (-2,0) -- (2,0) node[right] {$x$};
    \draw[->] (0,-2.5) -- (0,1) node[above] {$y$};

    % Horizontal line y=k
    \draw[red, thick] (-2,\k) -- (2,\k);
    \node[red, above left] at (-0.7,\k) {$y=k$, $k>0$};
\end{scope}

% Transformation arrow
\draw[->] (2.5, 0.5) -- (4.5, 0.5) node[midway, above] {$\frac{1}{z}$};

% W-plane
\begin{scope}[shift={(6,0)}]
    \draw[->] (-1.5,0) -- (1.5,0) node[right] {$u$};
    \draw[->] (0,-1) -- (0,2) node[above] {$v$};

    % Calculate center y and radius using pgfmath
    \tikzmath{
        \cy = -1/(2*\k);
        \radius = 1/(2*abs(\k)); % Use abs() for generality
    }

    % Circle centered at (0, cy) with radius r
    \draw[red, thick] (0, \cy) circle (\radius);

    % Center point
    \fill[red] (0, \cy) circle (1pt);

    % Radius arrow (pointing up-left as in image, angle approx 135 deg)
    \coordinate (CircPoint) at ({0 + \radius*cos(65)}, {\cy + \radius*sin(65)});
    \draw[red, thick, ->] (0, \cy) -- (CircPoint);

    % Radius label (position adjusted near the arrow)
    \node[red] at (-0.3, \cy+0.15) {$\frac{1}{2|k|}$};

    % Center label text (position adjusted to the right of the circle)
    \node[red, anchor=west, xshift=0.3cm] at (\radius, \cy) {center = $\left(0, \frac{1}{2k}\right)$};
\end{scope}

\end{tikzpicture}

\end{document}